\section*{Chapitre \ref{ch:m1:us-equity-market-structure} --- Structure du marché américain des actions}

\begin{solution}{exo:m1:us-equity-market-structure:1}
Le cours acheteur de D (80 actions) et le cours vendeur de C (50 actions) sont des \omterm{def:m1:us-equity-market-structure:lots}{lots irréguliers} et sont ignorés.
Meilleur cours acheteur 25.11 pour $100 + 300 = 400$ actions (B, C)~; meilleur cours vendeur 25.13 pour 300
actions (B)~: $25.11 \times 25.13$, 400 sur 300.
\end{solution}

\begin{solution}{exo:m1:us-equity-market-structure:2}
100 actions (4\,800~\$)~; 40 (10\,040~\$)~; 40 (39\,960~\$)~; 10 (42\,000~\$)~; 1
(712\,000~\$).
\end{solution}

\begin{solution}{exo:m1:us-equity-market-structure:3}
Bourses $49.4\% \times 17.6 = 8.7$ milliards d'actions~; dealers agissant en principal
$50.6\% \times 81.3\% \times 17.6 = 7.2$ milliards d'actions~; \omterm{def:m1:exchanges-brokers-venues:ats}{systèmes de négociation alternatifs}
1.7 milliard d'actions. Prix moyen $1.1\times10^{12}/17.6\times10^9 = \$62.50$.
\end{solution}

\begin{solution}{exo:m1:us-equity-market-structure:4}
Le cours vendeur protégé de B, à 25.13 pour 300 actions, a fait l'objet d'un \omterm{def:m1:us-equity-market-structure:opr}{trade-through}. Pour acheter 500 actions
d'un coup~: \omterm{def:m1:us-equity-market-structure:iso}{ordres de balayage intermarchés} (ISO) envoyés simultanément pour 300 à 25.13 chez B et
pour 200 à 25.14 chez A. (Les 50 actions de C à 25.13 ne sont pas protégées~; les prendre
aussi constitue une bonne exécution, et non une obligation.)
\end{solution}

\begin{solution}{exo:m1:us-equity-market-structure:5}
\omterm{def:m1:us-equity-market-structure:makertaker}{Tarification maker-taker}~: le preneur de liquidité paie 40.0030, le \omterm{def:m1:us-equity-market-structure:makertaker}{fournisseur de liquidité} reçoit 40.0025, la plateforme conserve 0.05
cent. \omterm{def:m1:us-equity-market-structure:makertaker}{Plateforme inversée}~: le preneur de liquidité paie 39.9990, le \omterm{def:m1:us-equity-market-structure:makertaker}{fournisseur de liquidité} reçoit 39.9986, la plateforme conserve 0.04
cent. Un \omterm{def:m1:us-equity-market-structure:makertaker}{fournisseur de liquidité} abandonne $0.25 + 0.14 = 0.39$ cent par action pour se faire exécuter en premier.
\end{solution}

\begin{solution}{exo:m1:us-equity-market-structure:6}
$\lambda\delta = 12 \times 600\times10^{-6} = 0.72\%$, soit 168 secondes par
jour. Lorsque les changements se regroupent, plusieurs tombent dans un même intervalle de longueur
$\delta$ et les intervalles périmés se chevauchent~: leur réunion est plus courte que leur
somme.
\end{solution}

\begin{solution}{exo:m1:us-equity-market-structure:7}
0.255~\% dans les deux cas. La péremption du flux consolidé par rapport à un \omterm{def:m1:us-equity-market-structure:sip}{flux direct} dépend
uniquement du délai supplémentaire $\delta$~: les deux vues attendent que les plateformes actualisent leurs cours.
Les délais des plateformes contrôlent autre chose --- combien de temps les \emph{cotations
elles-mêmes} retardent sur la fair value, ce qui définit la fenêtre pendant laquelle elles peuvent
faire l'objet de sniping.
\end{solution}

\begin{solution}{exo:m1:us-equity-market-structure:8}
(i) L'indice de référence est le NBBO du flux consolidé, périmé précisément quand les cours bougent~:
une exécution «~au NBBO~» pendant ces intervalles est moins bonne que le
vrai meilleur prix. (ii) L'horodatage à la réception ignore le temps que l'ordre passe ensuite
chez le wholesaler ou dans le routage~; la comparaison pertinente se fait avec le NBBO
au moment de l'exécution, et des écarts d'horloge d'une milliseconde suffisent à changer le résultat.
(iii) Pour les actions à cours élevé, les meilleurs cours étaient des \omterm{def:m1:us-equity-market-structure:lots}{lots irréguliers}, exclus du
NBBO~: l'indice de référence était plus large que le marché réel, si bien que «~au cours ou mieux~» était
facile. Plus généralement, «~au cours ou mieux~» ne dit rien ni de la taille ni du
point médian, ce que l'ordre aurait souvent pu obtenir.
\end{solution}

\begin{solution}{pb:m1:us-equity-market-structure:1}
\textbf{1.} $\lambda\delta = (8/60) \times 900\times10^{-6} = 0.012\%$.
\textbf{2.} 2.8 secondes.
\textbf{3.} La mise à jour de la plateforme la plus rapide, après 60 microsecondes~: un cours acheteur plus élevé
n'importe où, relève le meilleur cours acheteur.
\textbf{4.} Seulement lorsque la plateforme la \emph{plus lente} actualise son cours, après 500
microsecondes~: le meilleur cours vendeur est le plus bas, et les cours vendeurs bas périmés restent
les meilleurs jusqu'à ce que le dernier d'entre eux disparaisse.
\textbf{5.} Les quatre~: la plateforme la plus rapide actualise son cours à 60 microsecondes et le trader
arrive à 40 microsecondes.
\textbf{6.} Ancien cours vendeur $m+1$, nouvelle fair value $m+J$~: profit $J-1$ ticks, soit
0, 1 et 2 cents.
\textbf{7.} $0.25 \times 1 + 0.10 \times 2 = 0.45$ cent.
\textbf{8.} $0.45 \times 4 \times 300/100 = \$5.40$ par saut~; avec 3\,120
sauts par jour, 16\,848~\$.
\textbf{9.} $0.25\times0.70 + 0.10\times1.70 = 0.345$ cent~: 12\,917~\$ par jour.
Les sauts d'un tick ne couvrent pas les frais et sont laissés de côté~; seuls les sauts de deux ticks ou
plus sont exploités.
\textbf{10.} $0.45 \times 3 \times 3\,120 = \$4\,212$ par jour, 1.05 cent par
action négociée.
\textbf{11.} Non~: $0.35 - 1.05 = -0.70$ cent par action. Le \omterm{def:m1:what-a-trading-firm-does:market-maker}{teneur de marché} gagne 1\,400~\$ par jour
et perd 4\,212~\$.
\textbf{12.} Il économiserait 1.06 million de dollars par an, pour un coût de 2.5 millions de dollars~: pas
avec ces chiffres. Ses vrais choix sont de coter plus large, de coter en plus petites tailles, ou de
s'arrêter.
\textbf{13.} Chacun serait alors de nouveau plus lent que celui qui investira ensuite~; le
trader se modernise aussi, le profit du sniping est inchangé ou simplement réalloué,
et le secteur a dépensé 10 millions de dollars par an pour rien~: une course aux armements, dans
laquelle la vitesse n'a de valeur que relativement aux autres.
\textbf{14.} Une cotation est sûre si son \omterm{def:m1:what-a-trading-firm-does:market-maker}{teneur de marché} actualise son cours avant qu'un ordre retardé
puisse arriver, soit dans les $40 + 350 = 390$ microsecondes~: les trois
plus rapides. La plus lente, à 500 microsecondes, reste exposée.
\textbf{15.} Vers une plateforme qui affiche encore l'ancien cours vendeur, plus bas~; selon le cas,
l'ordre arrive après que la cotation a été actualisée ou prise et n'est pas exécuté, ou
s'exécute contre un cours vendeur plus élevé et récent. Le client court après le cours.
\textbf{16.} Non. La Rule 611 compare l'exécution aux cotations protégées
telles qu'affichées à l'instant~; le cours vendeur affiché de chaque plateforme était au niveau ou au-dessus du
prix payé. La règle protège les prix affichés, non les fair values, et ne dit
rien des cotations périmées.
\textbf{17.} $0.25\times0.90 + 0.10\times1.90 = 0.415$ cent~: un plafond de frais d'accès plus bas
rend le sniping \emph{plus} rentable.
\textbf{18.} Non~: l'opportunité vient de cotations en retard sur la fair value,
et non d'obligations de routage. L'abrogation changerait qui doit acheminer les ordres où,
et non la course à une cotation périmée.
\textbf{19.} \textbf{0.345 cent par action.}
\textbf{20.} Les investisseurs qui traitent avec les \omterm{def:m1:what-a-trading-firm-does:market-maker}{teneurs de marché}, à travers les
fourchettes plus larges et tailles plus petites qui financent les pertes de la question 10.
\end{solution}

\begin{solution}{iq:m1:us-equity-market-structure:1}
Le cours acheteur protégé le plus élevé et le cours vendeur protégé le plus bas sur l'ensemble des bourses, avec
les tailles. Officiellement, les SIP le calculent~; chaque société sérieuse calcule le sien
à partir des \omterm{def:m1:us-equity-market-structure:sip}{flux directs}. Les valeurs divergent parce que chaque observateur reçoit la mise à jour de chaque plateforme
après un délai différent et avec une horloge différente~: le NBBO est
défini par rapport à un lieu et à un instant, et «~le NBBO à 10:31:04.123456~»
n'a pas de réponse unique.
\iqlookfor{relativité de la vue consolidée, et non seulement «~le SIP est lent~».}
\end{solution}

\begin{solution}{iq:m1:us-equity-market-structure:2}
Un ordre à cours limité marqué pour indiquer à la plateforme destinataire d'exécuter sans
vérifier les marchés voisins, parce que l'émetteur a simultanément envoyé des ordres pour
balayer chaque cotation protégée de meilleur prix. Il existe pour qu'un ordre plus grand
que la meilleure taille affichée puisse consommer plusieurs niveaux de prix sur plusieurs plateformes
en un instant, au lieu d'être réacheminé niveau par niveau pendant que le marché
s'éloigne.
\iqlookfor{transfert de la responsabilité du respect des \omterm{def:m1:us-equity-market-structure:opr}{trade-through} vers
l'émetteur.}
\end{solution}

\begin{solution}{iq:m1:us-equity-market-structure:3}
Les preneurs de liquidité frappent d'abord la \omterm{def:m1:us-equity-market-structure:makertaker}{plateforme inversée} parce qu'elle les rémunère, si bien qu'un ordre au repos
y est exécuté avant le même prix sur n'importe quelle plateforme maker-taker~: la rémunération achète
la priorité temporelle entre plateformes. Elle vaut son coût lorsque la file d'attente au meilleur
prix est longue et que la fourchette est bloquée à un tick, de sorte que la priorité, et non
le prix, décide de qui négocie~--- et les exécutions obtenues tôt, avant que la file ait été
consommée, sont aussi celles qui subissent le moins la \omterm{def:m1:what-a-trading-firm-does:adverse-selection}{sélection adverse}.
\iqlookfor{la position dans la file comme ce qu'on achète.}
\end{solution}

\begin{solution}{iq:m1:us-equity-market-structure:4}
Par symbole, un petit tableau de taille fixe contenant les meilleures cotations par plateforme (prix sous forme
d'entiers), plus le meilleur cours acheteur et le meilleur cours vendeur en cache et des masques de bits par plateforme~; une mise à jour
ne touche qu'une case et, seulement si elle se trouvait au meilleur niveau ou l'améliore, un nouveau parcours de
seize entrées, qui tient dans une ou deux lignes de cache. Les symboles sont stockés dans un tableau dense
indexé par un identifiant entier, et non par un hachage de chaîne. Un thread par flux
décode et écrit~; on répartit les symboles entre les threads de construction des carnets d'ordres de sorte que chaque
symbole ait un seul écrivain et aucun verrou~; on publie vers les stratégies via un
instantané protégé par verrou de séquence ou un tampon circulaire. On épingle les threads, on alloue tout à l'avance,
on garde la structure chaude contiguë.
\iqlookfor{partitionnement à écrivain unique et disposition sensible au cache.}
\end{solution}

\begin{solution}{iq:m1:us-equity-market-structure:5}
Pour les investisseurs particuliers, en grande partie non~: ils obtiennent une petite amélioration de prix et aucune
commission. Pour les institutionnels, les \omterm{def:m1:us-equity-market-structure:dark}{dark pools} réduisent la fuite d'information. La
préoccupation porte sur la formation des prix~: les transactions hors bourse utilisent le NBBO sans
y contribuer, et si les ordres les moins informés sont écartés des bourses,
les fourchettes affichées s'élargissent pour tous ceux qui doivent y négocier, et
l'indice de référence de «~l'amélioration~» se dégrade. C'est aussi une question de concurrence~:
quelques wholesalers voient et tarifent l'essentiel des ordres de particuliers.
\iqlookfor{segmentation et \omterm{def:m1:what-a-trading-firm-does:adverse-selection}{sélection adverse}, dans les deux sens.}
\end{solution}

\begin{solution}{iq:m1:us-equity-market-structure:6}
(i) Les plateformes n'ont plus à réacheminer ni à rejeter les ordres qui entraîneraient un
\omterm{def:m1:us-equity-market-structure:opr}{trade-through}~: les petites bourses qui survivaient grâce au statut de cotation protégée perdent des flux,
et le volume se concentre~; positionner la connectivité et le capital de cotation sur les
plateformes les plus susceptibles de gagner. (ii) La meilleure exécution devient le jugement du courtier
plutôt qu'une règle mécanique~: les routeurs des courtiers et leurs analyses de coûts de transaction
comptent davantage, si bien qu'un \omterm{def:m1:what-a-trading-firm-does:market-maker}{teneur de marché} doit s'attendre à un contrôle plus strict de la qualité des exécutions
par plateforme et se démarquer sur ce terrain. (iii) Les marchés verrouillés et croisés deviennent
licites~: les prix affichés sur les différentes plateformes se chevaucheront brièvement, créant
un arbitrage au profit des plus rapides, ainsi que de nouveaux états que chaque moteur de cotation et chaque contrôle de risque
doit gérer. Il est également probable que se multiplient des modèles de plateformes plus différenciés (délais,
enchères par lots) désormais libres des exigences de cotation protégée.
\iqlookfor{effets de second ordre et actions concrètes des sociétés, et non une liste
d'opinions.}
\end{solution}
